\chapter{Маршрут из нескольких клеток}\label{ch:marshrut}

\textbf{Что вы узнете:} как пролететь несколько клеток подряд и нарисовать фигуру.

\section{Квадрат}

Чтобы нарисовать квадрат, дрону нужно побывать в четырёх клетках. Можно написать четыре \kod{code}, а можно одну команду \kod{code_s}: в ней клетки перечисляются через пробел в строке, а время делится между ними поровну.

\solfile{l07_square.py}

Квадрат 3$\times$3 клетки (рисунок~\ref{fig:quarter}): дрон взлетает в клетке 13, потом летит в клетки 16, 52, 49 и снова в 13. Каждая сторона 120\cm. Времени на четыре перехода 16 секунд, значит по 4 секунды на сторону, скорость 30\cm/с.

\begin{zametka}
Клетки в \kod{code_s} можно записывать списком \kod{[16, 52, 49, 13]} или строкой \kod{'16 52 49 13'}. Общее время делится поровну: на каждую клетку \kod{seconds / число клеток}.
\end{zametka}

\begin{vnimanie}[Направление]
Если смотреть сверху (клетка 0 внизу слева, строки растут вверх), маршрут 13, 16, 52, 49 идёт \emph{против} часовой стрелки: сначала вправо, потом вверх, потом влево.
\end{vnimanie}

\section{Как выбирать время}

\begin{itemize}
  \item Измерьте расстояние: считайте клетки (одна клетка 40\cm) или по диагонали — длиннее.
  \item Разделите на скорость. Для спокойного полёта берите 25–40\cm/с.
  \item Округлите вверх до целых секунд.
\end{itemize}

\begin{probuy}
  \item Треугольник: из клетки 13 в клетки 16, 52 и обратно в 13. Подберите время на каждую сторону. Сколько секунд длится шоу?
  \item Прямоугольник: 3 клетки в ширину, 2 в высоту (клетки 13, 16, 40, 37).
  \item Буква Z: сверху вправо (клетки 49, 52), по диагонали вниз (13), потом вправо (16).
  \item Квадрат в обратную сторону: 13, 49, 52, 16, 13.
\end{probuy}
\otvety{7}
